% === Chapter 0: Introduction ===
\section*{Introduction}
\label{sec:t0}
\addcontentsline{toc}{section}{Introduction}

The capabilities of Large Language Models (LLMs) and Vision-Language Models (VLMs)
are not fully unleashed during the pre-training stage alone.
Pre-training endows a model with statistical understanding of language and broad world knowledge,
but a model that has only undergone pre-training is essentially a high-quality next-token predictor ---
it does not know how to follow human instructions, cannot reliably distinguish harmful outputs from safe responses,
and is unable to exhibit structured reasoning processes in complex reasoning tasks.
\textbf{Post-training} is precisely the critical stage that bridges this gap:
it transforms a raw language model into a truly useful AI assistant.

\subsection{Definition and Importance of Post-Training}
\label{sec:t0-definition}

This book defines \textbf{post-training} as the collective term for all training stages
conducted after the completion of pre-training and before model deployment or inference,
aimed at enhancing specific model capabilities.
This definition encompasses the following core technical areas:

\begin{itemize}[leftmargin=2em]
  \item \textbf{Supervised Fine-Tuning (SFT)}:
    Using human-annotated (instruction, response) data to perform supervised fine-tuning on the model,
    enabling it to learn to follow instructions and produce well-formatted outputs
    \cite{wei2022flan,ouyang2022instructgpt,sanh2022multitask}.

  \item \textbf{Reinforcement Learning from Human Feedback (RLHF)}:
    Training a reward model to capture human preferences, then optimizing the model policy
    through reinforcement learning algorithms such as PPO,
    so that the outputs better align with human expectations
    \cite{christiano2017rlhf,ouyang2022instructgpt,bai2022constitutional}.

  \item \textbf{Direct Preference Optimization (DPO) and Its Variants}:
    Bypassing the explicit reward model and directly optimizing the policy from preference data pairs,
    substantially simplifying the training pipeline for alignment
    \cite{rafailov2023dpo,azar2023ipo,ethayarajh2024kto}.

  \item \textbf{Other Post-Training Techniques}:
    Including Constitutional AI \cite{bai2022constitutional},
    process reward models \cite{lightman2023prm},
    Group Relative Policy Optimization (GRPO) \cite{shao2024grpo},
    as well as specialized training methods targeting specific capabilities (reasoning, safety, tool use, multimodal).
\end{itemize}

\noindent
The boundaries between post-training and its upstream and downstream stages need to be clearly delineated:

\begin{itemize}[leftmargin=2em]
  \item \textbf{Distinction from Pre-Training}:
    Pre-training is conducted on large-scale unlabeled corpora, with the objective of general language modeling (typically next-token prediction or masked language modeling).
    Post-training, in contrast, uses relatively small-scale but high-quality annotated data or preference signals, with the objective of endowing the model with specific behavioral patterns.
    The two differ by orders of magnitude in data scale, training objectives, and computational cost.

  \item \textbf{Distinction from Inference-Time Techniques}:
    Techniques such as chain-of-thought prompting \cite{wei2022cot},
    in-context learning \cite{brown2020gpt3},
    and retrieval-augmented generation (RAG) \cite{lewis2020rag}
    improve model performance at inference time through prompt engineering or external knowledge retrieval,
    but do not modify model parameters. Post-training, by contrast, permanently alters model weights through gradient updates.
    Of course, the two are not entirely orthogonal ---
    for example, OpenAI o1 \cite{openai2024o1} deeply integrates the reasoning capabilities acquired during post-training
    with inference-time extended compute strategies.
\end{itemize}

\noindent
The importance of post-training can be understood from three dimensions.
\textbf{First}, it serves as the bridge from ``capability'' to ``usability'' ---
InstructGPT \cite{ouyang2022instructgpt}, with a mere 1.3B-parameter SFT+RLHF model, outperformed
the 175B original GPT-3 in human evaluations, directly demonstrating that post-training can achieve
qualitative leaps without increasing model scale.
\textbf{Second}, it is the core mechanism for implementing alignment ---
from safety alignment to value calibration, from reducing hallucinations to improving reasoning reliability,
post-training is the key stage that translates abstract alignment objectives into measurable training signals.
\textbf{Third}, it is the primary battleground for model differentiation ---
as pre-training architectures and data converge toward homogeneity (predominantly decoder-only Transformers + web-scale corpora),
differences in post-training strategies increasingly become the decisive factor in determining final product quality.

\paragraph{Target Audience of This Book.}
This book is intended for three types of readers:
(1)~\textbf{ML researchers}, who seek a systematic understanding of the latest advances across post-training subfields and their interconnections;
(2)~\textbf{graduate students}, who need a structured learning path from foundational concepts to frontier methods;
(3)~\textbf{industry practitioners}, who are concerned with practical selection and deployment strategies for post-training techniques in model development.
This book assumes that readers possess foundational knowledge of Transformer architectures and language modeling,
but does not require prior familiarity with each subfield (e.g., RLHF, multimodal training, agent tuning).

\paragraph{Distinction from Existing Surveys.}
Several representative surveys related to post-training already exist.
\citet{casper2023rlhfproblems} systematically cataloged the open problems and fundamental limitations of RLHF,
but the scope was limited to reward modeling and preference learning,
without addressing subsequent branches such as multimodal and agentic post-training.
\citet{wang2023alignsurvey} organized LLM alignment techniques along three dimensions --- data collection, training methods, and evaluation ---
providing a clear taxonomy, but adopted a static classification rather than an evolutionary perspective.
\citet{shen2023alignmentsurvey} approached the topic from the theoretical framework of inner/outer alignment,
covering interpretability and adversarial attacks, but did not delve into the technical details of individual subfields.
This book complements the above works in three respects:
(1)~\textbf{broader coverage}, integrating SFT, preference optimization, safety, reasoning, multimodal, agentic, and efficiency into a unified framework across seven threads;
(2)~an \textbf{evolutionary structure} rather than a static taxonomy, tracing the temporal evolution of each thread from problem origin to the frontier;
(3)~\textbf{cross-thread references} that explicitly document the flow of ideas across subfields, revealing technical transfers that traditional classification approaches tend to obscure.

\subsection{Organization of This Book}
\label{sec:t0-organization}

Traditional surveys are typically organized by method (e.g., one chapter for SFT, one for RLHF, one for DPO).
This organizational approach is clear but static, and fails to capture the most fundamental characteristic of the post-training landscape:
\textbf{dense idea flow and technical transfer across different research directions}.

For example, the ideas behind reward modeling flow from preference optimization (Thread~2)
to process reward models in reasoning verification (Thread~4),
and then loop back to preference optimization through GRPO;
the instruction tuning pipeline developed in SFT (Thread~1)
is directly adopted and extended by multimodal post-training (Thread~5) and tool use (Thread~6);
efficiency techniques such as LoRA (Thread~7) enable the methods across all preceding threads to be deployed at lower cost.

To capture this dynamic structure, this book adopts a \textbf{Hybrid Problem-Evolution Graph} as its organizational framework.
The core ideas are as follows:

\begin{enumerate}[leftmargin=2em]
  \item \textbf{Seven Problem-Driven Evolution Threads}:
    This book divides the post-training landscape into seven relatively independent problem threads,
    each centered around a core question.
    Within each thread, the evolution is traced chronologically along the chain of
    ``Problem Origin $\rightarrow$ Foundation $\rightarrow$ Limitations \& Branching $\rightarrow$ Convergence $\rightarrow$ Open Problems''.

  \item \textbf{Cross-Thread References}:
    When a technique or idea from one thread is adopted by another,
    this book annotates it using \crossref{1}{sec:t1}-style cross-references,
    enabling readers to trace the origin and destination of ideas.

  \item \textbf{Panoramic View}:
    \cref{fig:overview} presents a panoramic view of the seven threads as a TikZ diagram,
    where horizontal swim lanes correspond to threads, the x-axis represents the timeline,
    nodes denote landmark works, and dashed arrows indicate cross-thread idea flow.
\end{enumerate}

\noindent
The advantages of this organizational approach are threefold:
(1)~it maintains a clear timeline and causal relationships within each thread;
(2)~through cross-thread references, it reveals intellectual connections that traditional taxonomies tend to sever;
(3)~the panoramic diagram provides readers with an entry point to begin reading from any position.

\subsection{Overview of the Seven Evolution Threads}
\label{sec:t0-threads}

\paragraph{Thread 1: Instruction Following \& SFT.}
Core question: \textit{How can a pre-trained language model be taught to understand and follow human instructions?}
FLAN \cite{wei2022flan} and T0 \cite{sanh2022multitask} first demonstrated that
multi-task instruction tuning can significantly improve a model's zero-shot generalization capability.
InstructGPT \cite{ouyang2022instructgpt} established SFT as the first stage of the RLHF pipeline,
codifying the classic three-stage training paradigm of ``SFT $\rightarrow$ reward model $\rightarrow$ PPO''.
Self-Instruct \cite{wang2023selfinstruct} pioneered the method of leveraging the model itself to generate instruction data,
drastically reducing annotation costs and giving rise to open-source instruction-tuned models such as Alpaca \cite{taori2023alpaca}.
LIMA \cite{zhou2023lima} achieved highly competitive performance with only 1,000 carefully curated samples,
proposing the ``Superficial Alignment Hypothesis'' --- that the knowledge required for alignment has already been acquired during pre-training,
and the role of SFT is merely to activate the correct output format and style.
\crossref{1}{sec:t1}

\paragraph{Thread 2: Reward Modeling \& Preference Optimization.}
Core question: \textit{How can vague human preferences be transformed into optimizable training signals?}
The foundational framework of RLHF can be traced back to
the work of Christiano et al.\ \cite{christiano2017rlhf} on Atari and MuJoCo tasks,
which was subsequently applied to text summarization by Stiennon et al.\ \cite{stiennon2020summarize},
and comprehensively applied to dialogue model training in InstructGPT \cite{ouyang2022instructgpt}.
However, the RLHF training pipeline is complex and unstable, involving the coordinated training of four separate models.
DPO \cite{rafailov2023dpo} demonstrated through mathematical derivation that the reward model can be implicitly parameterized,
thereby converting RLHF into a simple classification loss and substantially lowering the implementation barrier.
Subsequently, variants such as IPO \cite{azar2023ipo}, KTO \cite{ethayarajh2024kto},
and ORPO \cite{hong2024orpo} addressed various limitations of DPO from different angles.
GRPO \cite{shao2024grpo} introduced group-relative rewards into policy optimization,
achieving breakthrough results on reasoning tasks.
In 2025, the emergence of the RLVR (Reinforcement Learning with Verifiable Rewards) paradigm
marked a new watershed: methods such as DAPO \cite{yu2025dapo} entirely dispensed with reward models,
directly using rule-based verifiers as reward signals, improving GRPO's performance on AIME from 30\% to 50\%.
\crossref{2}{sec:t2}

\paragraph{Thread 3: Safety \& Alignment.}
Core question: \textit{How can a model be ensured to refuse harmful requests while following instructions, without becoming overly conservative?}
Constitutional AI (CAI) \cite{bai2022constitutional} proposed having the model critique and revise its own outputs
according to a set of predefined principles (a constitution),
reducing the dependence of safety alignment on human annotation of harmful data.
Safe RLHF \cite{dai2024saferlhf} decoupled helpfulness and harmlessness into two independent objectives,
seeking an optimal balance on the Pareto frontier through a constrained optimization framework.
Meanwhile, the evolution of jailbreaking attacks ---
from simple prompt injection to gradient-based automated attacks such as GCG \cite{zou2023gcg} ---
has created an ongoing adversarial arms race, driving continuous iteration in safety research.
In 2025--2026, the emergence of reasoning models introduced a new safety dimension:
research has shown that reasoning models can be employed as autonomous jailbreak agents,
achieving a 97\% attack success rate with zero human intervention \cite{hagendorff2026jailbreakagents},
making the safety of chain-of-thought reasoning a new core challenge.
\crossref{3}{sec:t3}

\paragraph{Thread 4: Reasoning \& Verification.}
Core question: \textit{How can models perform reliable multi-step reasoning and verify intermediate steps?}
Chain-of-Thought (CoT) prompting \cite{wei2022cot} demonstrated that
including reasoning processes in the prompt can significantly improve model performance on mathematical and logical tasks,
but the upper bound of this capability is constrained by inference-time prompt design.
Process Reward Models (PRMs) \cite{lightman2023prm} refined the verification granularity from final answers to each individual reasoning step,
providing more fine-grained signals for reward-guided search.
OpenAI o1 \cite{openai2024o1} deeply integrated reasoning capabilities into post-training,
training models through large-scale reinforcement learning to autonomously unfold long chains of reasoning at inference time,
achieving expert-level performance on mathematics competitions and programming tasks.
DeepSeek-R1 \cite{deepseek2025r1} further explored
the possibility of training reasoning capabilities through pure RL (without SFT cold start),
revealing the remarkable phenomenon that reasoning capabilities can spontaneously emerge from RL.
In 2025--2026, thinking models became the default paradigm:
OpenAI o3 \cite{openai2025o3} achieved 88.9\% on AIME 2025,
DeepSeek-V3.2 \cite{deepseek2025v32} reached gold-medal level on IMO-level mathematics competitions,
and Qwen3 \cite{yang2025qwen3} achieved seamless integration of thinking and non-thinking modes.
\crossref{4}{sec:t4}

\paragraph{Thread 5: Multimodal Post-training.}
Core question: \textit{How can text-only post-training techniques be extended to joint vision-language modeling?}
LLaVA \cite{liu2023llava} introduced visual instruction tuning into VLM training,
constructing an efficient multimodal instruction-following model
through a two-stage training process (visual-text alignment + instruction tuning).
Subsequently, RLHF techniques were introduced into VLMs to reduce multimodal hallucination ---
i.e., the model generating descriptions inconsistent with the image content
\cite{sun2023rlhfv,yu2024rlhfv}.
Hallucination reduction has become one of the core challenges of multimodal post-training,
giving rise to dedicated evaluation benchmarks such as POPE \cite{li2023pope}
and CHAIR, as well as targeted training methods.
In 2025, RL for VLM reasoning emerged as a new frontier direction:
works such as R1-VL \cite{zhang2025r1vl} and MM-Eureka \cite{meng2025mmeureka}
extended the RLVR paradigm from text-only reasoning to multimodal reasoning,
marking the paradigm shift in VLM post-training from SFT+DPO to SFT+RL.
\crossref{5}{sec:t5}

\paragraph{Thread 6: Tool Use \& Agentic Capabilities.}
Core question: \textit{How can post-training enable models to invoke external tools and complete complex tasks?}
Toolformer \cite{schick2023toolformer} demonstrated that models can learn through self-supervised learning
to insert API calls during the generation process.
The function calling mechanism introduced by OpenAI in GPT-3.5/4
standardized tool use as a structured JSON schema interface,
making it a foundational infrastructure for production-grade applications.
Agent tuning \cite{zeng2023agenttuning} further extended the training objective from single tool invocations
to multi-step task planning and execution, laying the groundwork for building autonomous agents.
This thread has deep intersections with Thread~1 (instruction following provides the foundational capability)
and Thread~4 (reasoning provides the planning capability).
In 2024--2025, agent capabilities evolved from API calls to direct manipulation of computer GUIs:
Claude Computer Use \cite{anthropic2024computeruse} and OpenAI Operator \cite{openai2025operator}
demonstrated the ability of AI agents to complete complex tasks in real desktop environments,
and standardized protocols such as MCP \cite{anthropic2024mcp} are driving agents from experimental prototypes toward industrial deployment.
\crossref{6}{sec:t6}

\paragraph{Thread 7: Efficiency \& Data Engineering.}
Core question: \textit{How can equivalent or superior post-training results be achieved with lower computational cost and less data?}
LoRA \cite{hu2022lora} reduced the number of trainable parameters for fine-tuning to 0.1\%--1\% of the original model through low-rank decomposition,
making post-training feasible on consumer-grade GPUs.
QLoRA \cite{dettmers2023qlora} further combined quantization techniques,
enabling fine-tuning of 65B models within 48GB of GPU memory.
On the data engineering front,
data selection methods (such as DEITA \cite{liu2024deita}
and Instruction Mining \cite{cao2024instructionmining}) aim to identify the most valuable subsets from large-scale instruction data;
synthetic data generation leverages strong models (such as GPT-4) to produce training data,
giving rise to methodologies such as Alpaca \cite{taori2023alpaca},
WizardLM \cite{xu2024wizardlm},
and Evol-Instruct.
\crossref{7}{sec:t7}

\subsection{Cross-Thread References}
\label{sec:t0-crossref}

The seven evolution threads do not develop in isolation; dense idea flow and technical transfer exist among them.
\cref{fig:overview} presents these cross-thread connections in a panoramic view.
Here we outline the most important idea flows:

\begin{itemize}[leftmargin=2em]
  \item \textbf{T1 $\rightarrow$ T2} (SFT provides a warm start for RLHF):
    The ``SFT $\rightarrow$ RM $\rightarrow$ PPO'' pipeline established by InstructGPT
    positions instruction tuning as the preceding stage of preference optimization,
    a paradigm inherited by virtually all subsequent works.

  \item \textbf{T1 $\rightarrow$ T3} (Instruction following exposes safety risks):
    The instruction-following capability conferred by SFT in T1 simultaneously opens safety vulnerabilities ---
    a model that faithfully follows instructions will also faithfully execute malicious ones.
    This inherent tension is precisely the problem origin of T3,
    driving the development of safety alignment methods such as Constitutional AI.

  \item \textbf{T2 $\rightarrow$ T3} (Preference optimization shapes safety):
    RLHF and DPO provide a general framework for converting human safety preferences into training signals,
    upon which Constitutional AI and Safe RLHF developed specialized methods for safety.

  \item \textbf{T2 $\leftrightarrow$ T4} (Bidirectional flow between reward modeling and reasoning):
    The ideas behind reward models flow from T2 into T4, giving rise to process reward models;
    conversely, the verifiable rewards developed in T4 (automatic verification based on ground-truth answers)
    inspired methods such as GRPO in T2, enabling preference optimization without human annotation.

  \item \textbf{T1 $\rightarrow$ T5, T6} (Cross-modal transfer of instruction tuning):
    The mature instruction tuning pipeline developed in T1 was directly adopted by T5 (LLaVA's visual instruction tuning)
    and T6 (Toolformer's tool use instruction tuning),
    demonstrating the high transferability of this paradigm.

  \item \textbf{T7 $\rightarrow$ T1--T6} (Democratizing effect of efficiency techniques):
    Parameter-efficient training methods such as LoRA and QLoRA have enabled the post-training techniques
    across the preceding six threads to be widely reproduced and iterated upon in academia and the open-source community,
    significantly accelerating progress across the entire field.

  \item \textbf{T4 $\rightarrow$ T6} (Reasoning empowers agent planning):
    The long-chain reasoning capabilities developed in T4 provide core support for multi-step agent planning in T6;
    the reasoning capabilities of o1 and DeepSeek-R1 directly enhance agent performance on complex tasks.

  \item \textbf{T2 $\rightarrow$ T4 $\rightarrow$ T5} (Cross-thread migration of RLVR):
    In 2025, the RLVR paradigm was systematized in T2 (DAPO \cite{yu2025dapo}),
    applied at scale to reasoning training in T4,
    and subsequently migrated rapidly to T5 (R1-VL \cite{zhang2025r1vl}, MM-Eureka \cite{meng2025mmeureka}),
    forming a clear T2 $\rightarrow$ T4 $\rightarrow$ T5 technical migration path.
\end{itemize}

% === Overview: Hybrid Problem-Evolution Graph ===
% 7 horizontal lanes (one per thread), shared x-axis 2017--2026,
% landmark papers as nodes, dashed cross-thread arrows.

\begin{figure}[htbp]
\centering
\resizebox{\textwidth}{!}{%
\begin{tikzpicture}[
    >=Stealth,
    node distance=0.6cm,
    paper/.style={
        rounded corners=2pt,
        font=\footnotesize\bfseries,
        inner sep=4pt,
        minimum height=0.65cm,
        text=white,
        drop shadow={shadow xshift=0.5pt, shadow yshift=-0.5pt, opacity=0.3},
    },
    crossarrow/.style={
        ->,
        dashed,
        very thick,
        shorten >=2pt,
        shorten <=2pt,
    },
    arrowlabel/.style={
        font=\tiny\itshape,
        fill=white,
        fill opacity=0.85,
        text opacity=1,
        inner sep=1pt,
        rounded corners=1pt,
    },
    lanelabel/.style={
        font=\small\bfseries,
        anchor=east,
        text width=4.2cm,
        align=right,
    },
]

% --- Geometry ---
\def\laneheight{1.7}   % vertical spacing between lanes
\def\xmin{0}
\def\xmax{20}          % total width
% Year mapping: 2017 = 0, 2026 = 20  =>  1 year = 2 units
\newcommand{\yearx}[1]{\fpeval{(#1-2017)*2}}

% --- Year axis ---
\foreach \y in {2017,2018,2019,2020,2021,2022,2023,2024,2025,2026} {
    \draw[gray!40] (\yearx{\y}, 0.5*\laneheight) -- (\yearx{\y}, -6.5*\laneheight);
    \node[font=\footnotesize, gray!80, anchor=south] at (\yearx{\y}, 0.7*\laneheight) {\y};
}

% --- Lane backgrounds ---
\foreach \i/\col in {0/thread1,1/thread2,2/thread3,3/thread4,4/thread5,5/thread6,6/thread7} {
    \fill[\col, opacity=0.06]
        (\xmin, -\i*\laneheight+0.45*\laneheight)
        rectangle
        (\xmax, -\i*\laneheight-0.45*\laneheight);
}

% =====================================================================
% Thread 1: Instruction Following & SFT (lane y=0)
% =====================================================================
\def\laneA{0}
\node[lanelabel, text=thread1] at (-0.3, \laneA*\laneheight)
    {T1: Instruction\\Following \& SFT};

\node[paper, fill=thread1] (flan)
    at (\yearx{2021}+0.5, \laneA*\laneheight) {FLAN};
\node[paper, fill=thread1] (instructgpt-sft)
    at (\yearx{2022}+0.8, \laneA*\laneheight) {InstructGPT};
\node[paper, fill=thread1] (selfinstruct)
    at (\yearx{2023}+1.2, \laneA*\laneheight) {Self-Inst.};
\node[paper, fill=thread1] (lima)
    at (\yearx{2024}+0.5, \laneA*\laneheight) {LIMA};

% =====================================================================
% Thread 2: Reward Modeling & Preference Optimization (lane y=-1)
% =====================================================================
\def\laneB{-1}
\node[lanelabel, text=thread2] at (-0.3, \laneB*\laneheight)
    {T2: Reward Modeling \&\\Preference Opt.};

\node[paper, fill=thread2] (rlhf-og)
    at (\yearx{2017}+1.0, \laneB*\laneheight) {RLHF};
\node[paper, fill=thread2] (instructgpt-rlhf)
    at (\yearx{2022}+0.8, \laneB*\laneheight) {InstructGPT};
\node[paper, fill=thread2] (dpo)
    at (\yearx{2023}+0.6, \laneB*\laneheight) {DPO};
\node[paper, fill=thread2] (grpo)
    at (\yearx{2024}+1.0, \laneB*\laneheight) {GRPO};
\node[paper, fill=thread2] (dapo)
    at (\yearx{2025}+1.0, \laneB*\laneheight) {DAPO/RLVR};

% =====================================================================
% Thread 3: Safety & Alignment (lane y=-2)
% =====================================================================
\def\laneC{-2}
\node[lanelabel, text=thread3] at (-0.3, \laneC*\laneheight)
    {T3: Safety \&\\Alignment};

\node[paper, fill=thread3] (cai)
    at (\yearx{2022}+0.8, \laneC*\laneheight) {Const.\ AI};
\node[paper, fill=thread3] (saferlhf)
    at (\yearx{2023}+1.2, \laneC*\laneheight) {Safe RLHF};
\node[paper, fill=thread3] (jailbreak)
    at (\yearx{2024}+1.2, \laneC*\laneheight) {Jailbreak};

% =====================================================================
% Thread 4: Reasoning & Verification (lane y=-3)
% =====================================================================
\def\laneD{-3}
\node[lanelabel, text=thread4] at (-0.3, \laneD*\laneheight)
    {T4: Reasoning \&\\Verification};

\node[paper, fill=thread4] (cot)
    at (\yearx{2022}+0.3, \laneD*\laneheight) {CoT};
\node[paper, fill=thread4] (prm)
    at (\yearx{2023}+0.5, \laneD*\laneheight) {PRM};
\node[paper, fill=thread4] (oone)
    at (\yearx{2024}+0.0, \laneD*\laneheight) {o1};
\node[paper, fill=thread4] (deepseekr1)
    at (\yearx{2025}-0.3, \laneD*\laneheight) {R1};
\node[paper, fill=thread4] (othree)
    at (\yearx{2025}+0.8, \laneD*\laneheight) {o3};
\node[paper, fill=thread4] (v32)
    at (\yearx{2026}+0.0, \laneD*\laneheight) {V3.2};

% =====================================================================
% Thread 5: Multimodal Post-training (lane y=-4)
% =====================================================================
\def\laneE{-4}
\node[lanelabel, text=thread5] at (-0.3, \laneE*\laneheight)
    {T5: Multimodal\\Post-training};

\node[paper, fill=thread5] (llava)
    at (\yearx{2023}-0.3, \laneE*\laneheight) {LLaVA};
\node[paper, fill=thread5] (vlmrlhf)
    at (\yearx{2024}+0.0, \laneE*\laneheight) {VLM-RLHF};
\node[paper, fill=thread5] (halluc)
    at (\yearx{2025}+0.3, \laneE*\laneheight) {Anti-Halluc.};

% =====================================================================
% Thread 6: Tool Use & Agentic (lane y=-5)
% =====================================================================
\def\laneF{-5}
\node[lanelabel, text=thread6] at (-0.3, \laneF*\laneheight)
    {T6: Tool Use \&\\Agentic};

\node[paper, fill=thread6] (toolformer)
    at (\yearx{2022}+1.5, \laneF*\laneheight) {Toolformer};
\node[paper, fill=thread6] (funccall)
    at (\yearx{2023}+1.5, \laneF*\laneheight) {Func.\ Call};
\node[paper, fill=thread6] (agenttune)
    at (\yearx{2024}+1.5, \laneF*\laneheight) {AgentTune};
\node[paper, fill=thread6] (mcp)
    at (\yearx{2025}+1.3, \laneF*\laneheight) {MCP/Agents};

% =====================================================================
% Thread 7: Efficiency & Data Engineering (lane y=-6)
% =====================================================================
\def\laneG{-6}
\node[lanelabel, text=thread7] at (-0.3, \laneG*\laneheight)
    {T7: Efficiency \&\\Data Engineering};

\node[paper, fill=thread7] (lora)
    at (\yearx{2021}+1.0, \laneG*\laneheight) {LoRA};
\node[paper, fill=thread7] (datasel)
    at (\yearx{2023}+0.8, \laneG*\laneheight) {Data Selection};
\node[paper, fill=thread7] (synthdata)
    at (\yearx{2024}+0.5, \laneG*\laneheight) {Synthetic Data};

% =====================================================================
% Cross-thread arrows (dashed, idea flow)
% =====================================================================

% T1 -> T2: InstructGPT-SFT feeds into InstructGPT-RLHF
\draw[crossarrow, thread1] (instructgpt-sft.south) -- (instructgpt-rlhf.north);

% T2 -> T3: RLHF ideas feed into Constitutional AI
\draw[crossarrow, thread2] (instructgpt-rlhf.south) -- (cai.north)
    node[arrowlabel, midway, right=1pt] {reward};

% T2 -> T4: Reward modeling feeds reasoning verification (PRM)
\draw[crossarrow, thread2] (dpo.south)
    to[out=-90, in=90] (prm.north);

% T4 -> T2: Reasoning verification inspires GRPO
\draw[crossarrow, thread4] (prm.east)
    to[out=0, in=180] (grpo.south west)
    node[arrowlabel, midway, above=1pt] {verification};

% T1 -> T5: SFT pipeline adopted by LLaVA
\draw[crossarrow, thread1] (selfinstruct.south)
    to[out=-90, in=90] (llava.north)
    node[arrowlabel, near start, right=1pt] {SFT pipeline};

% T2 -> T5: RLHF adopted for VLM alignment
\draw[crossarrow, thread2] (dpo.south)
    to[out=-100, in=90] (vlmrlhf.north)
    node[arrowlabel, near end, left=1pt] {preference};

% T1 -> T6: Instruction tuning enables tool use
\draw[crossarrow, thread1] (selfinstruct.south)
    to[out=-90, in=90] (toolformer.north);

% T7 -> T1: LoRA enables efficient SFT
\draw[crossarrow, thread7] (lora.north)
    to[out=90, in=-90] (flan.south)
    node[arrowlabel, near start, right=1pt] {LoRA};

% T7 -> T2: Efficient training enables broader RLHF
\draw[crossarrow, thread7] (datasel.north)
    to[out=90, in=-90] (grpo.south);

% T3 -> T4: Safety constraints shape reasoning alignment
\draw[crossarrow, thread3] (saferlhf.south)
    to[out=-90, in=90] (oone.north)
    node[arrowlabel, midway, right=1pt] {safety};

% T4 -> T6: Reasoning enables agent planning
\draw[crossarrow, thread4] (oone.south)
    to[out=-90, in=90] (agenttune.north)
    node[arrowlabel, midway, right=1pt] {reasoning};

% T2 -> T4: RLVR (DAPO) feeds reasoning training
\draw[crossarrow, thread2] (dapo.south)
    to[out=-90, in=90] (othree.north)
    node[arrowlabel, near start, right=1pt] {RLVR};

% T4 -> T5: Reasoning RL migrates to multimodal (R1-VL, MM-Eureka)
\draw[crossarrow, thread4] (deepseekr1.south)
    to[out=-90, in=90] (halluc.north east);

% T4 -> T6: Thinking models empower agent planning
\draw[crossarrow, thread4] (othree.south)
    to[out=-90, in=90] (mcp.north);

\end{tikzpicture}%
}% end resizebox
\caption{Post-training 领域的 Hybrid Problem-Evolution Graph 全景视图。
七条水平泳道分别对应七条演化线程，
x 轴为时间线（2017--2026），
节点为各线程内的 landmark 工作，
虚线箭头表示跨线程的思想流动与技术迁移。}
\label{fig:overview}
\end{figure}


\noindent
In the chapters that follow, we will conduct in-depth analyses of the technical details of landmark papers along each thread,
while maintaining a global perspective through cross-thread references. Readers may read linearly following the thread order (T1--T7),
or start from any thread of interest and navigate to related threads via cross-references.
